\documentclass[10pt,a4paper]{amsart}
\usepackage[margin=19mm]{geometry}
\usepackage{amsmath,amssymb,mathtools}
\usepackage[scr=rsfs,cal=euler]{mathalfa}
\usepackage{xfrac}
\usepackage[unicode,hidelinks,bookmarks=false]{hyperref}
\newcommand{\ie}{i.e.\ }
\newcommand{\cf}{cf.\ }
\newcommand{\resp}{respectively\ }
\newcommand{\Subsection}{\subsection}
\providecommand{\nobreakdash}{\nobreak\hbox{-}}
\setlength{\emergencystretch}{2em}
\multlinegap0pt
\usepackage{bm}
\usepackage{bbm}
\usepackage{dsfont}
\usepackage{xspace}
\usepackage{cancel}
\usepackage{mathtools}
\usepackage{amsthm}
\makeatletter
\@ifundefined{lemma}{\newtheorem{lemma}{Lemma}}{}
\@ifundefined{proposition}{\newtheorem{proposition}{Proposition}}{}
\makeatother
\title[Adjoint-Compatible Surrogates of the Expected Information Ga]{Adjoint-Compatible Surrogates of the Expected Information Gain for Optimal Experimental Design}
\author{Luc de Montella, Sebastian Sager}
\date{Source: arXiv:2603.26431v3 (CC BY 4.0)}
\begin{document}
\maketitle

\noindent\emph{Source and licence.} Adjoint-Compatible Surrogates of the Expected Information Gain for Optimal Experimental Design. Version arXiv:2603.26431v3 (CC BY 4.0), distributed under the Creative Commons Attribution 4.0 International licence, as recorded in the arXiv licence metadata for this exact version and shown on the abstract page https://arxiv.org/abs/2603.26431v3. Full rights notice: licenses/arxiv-2603.26431-CC-BY-4.0.txt.\par\smallskip

\noindent\textit{Editorial scope.} Counted unit: the Proposition whose source label is \texttt{prop:consistency\_LG} in the article (source0.tex lines 879--908, source label \texttt{prop:consistency\_LG}), delivered together with the article's own complete original proof of it (source0.tex lines 909--1036, one \texttt{proof} environment of 128 lines). No mathematics is written, repaired or re-derived: statement and proof are the article's own text line for line. Required context retained uncounted: eq:accumulator\_definition (lines 688--692); lmm:weak\_conv\_qk (lines 1985--2032); lmm:behavior\_denom\_eig\_gauss (lines 2098--2129). Every definition, display and label of those lines is retained unchanged. Omitted from this package and not redistributed: 1--687, 693--878, 1037--1984, 2033--2097, 2130--2242. Nothing inside a retained excerpt is omitted or altered: every retained body is delivered exactly as the article's lines are written, so no omission receipt of any kind accompanies this package. Citations: the retained excerpts carry no citation command at all, so no bibliography entry of the article is transcribed and no reference is redistributed. Reference adaptation: none was needed and none was made; every reference inside the delivered unit resolves inside it, so no unresolved reference or citation is produced. Printed slips kept literal: no printed word, symbol, hypothesis or number of the counted statement or of its proof was altered. Layout and class: the article's own class is \texttt{article} with packages including fontenc, inputenc, lmodern, bm, geometry, amsthm, mathtools, biblatex, amsmath, amssymb, amsfonts, hyperref, xparse, booktabs; the wrapper is the lane's pinned \texttt{amsart} editorial wrapper at 10pt on A4, which restates the theorem-like environments the retained text needs and adds the article's own macro definitions that the retained text needs (none), with extra wrapper packages (bm, bbm, dsfont, xspace, cancel, mathtools). The delivered numbering is the unit's own. Assets: no retained span contains a figure environment, a graphics-inclusion command, a PDF-page inclusion, a table or a TikZ picture; no image file is copied into this package, the package declares no asset dependency and no image file is redistributed with it. Licence: version arXiv:2603.26431v3 of this article is distributed under the Creative Commons Attribution 4.0 International licence, as recorded in the arXiv licence metadata for this exact version and shown on the abstract page https://arxiv.org/abs/2603.26431v3. A full rights notice is delivered at licenses/arxiv-2603.26431-CC-BY-4.0.txt. Characters: every retained excerpt is pure ASCII, so the delivered engine, XeLaTeX with its default Latin Modern text font, reports no missing character.
\begin{equation} \label{eq:accumulator_definition}
F_0 = 0,
\qquad
F_{i+1} = F_{i} + F_{i+1}^\Delta,
\end{equation}

\begin{lemma} \label{lmm:weak_conv_qk}
Let $(p_k)_{k\ge1}$ be a sequence of probability measures on $\mathbb{R}^d$  with mean $m_k$ and finite second moment such that
\[
p_k \to p_0 := \mathcal N(m_0,\Sigma_0)
\quad\text{in } W_2.
\]

Let $F$ be a symmetric positive semidefinite matrix, and consider the probability measures 

\[
q^{(k)}(d\theta)
=
\frac{\phi^{(k)}(\theta)}{Z^{(k)}}\,p_k(d\theta)\quad\text{and}\quad
q^{(0)}(d\theta)
=
\frac{\phi(\theta)}{Z^{(0)}}\,p_0(d\theta)
\]


where
\[
\phi^{(k)}(\theta)
=
\exp\!\left(-\tfrac12(\theta-m_k)^\top F (\theta-m_k)\right),
\qquad
\phi(\theta)
=
\exp\!\left(-\tfrac12(\theta-m_0)^\top F (\theta-m_0)\right),
\]
and
\[
Z^{(k)}=\int \phi^{(k)}(\theta)p_k(d\theta),
\qquad
Z^{(0)}=\int \phi(\theta)p_0(d\theta).
\]

Let $(f_k)$ be continuous functions converging locally uniformly to $f$ and satisfying
\[
|f_k(\theta)| \le C(1+\|\theta\|^2).
\]

Then
\[
\mathbb{E}_{q^{(k)}}[f_k(\theta)]
\;\longrightarrow\;
\mathbb{E}_{q^{(0)}}[f(\theta)] .
\]
\end{lemma}

\begin{lemma}\label{lmm:behavior_denom_eig_gauss}
Let $(p_k)_{k\ge1}$, $q^{(k)}$, and $q^{(0)}$ be as in Lemma~\ref{lmm:weak_conv_qk}.
Let \(\ell(y\mid\theta)\) be the likelihood associated with the linear-Gaussian observation model
\[
y_d = H_{d}\theta + b_{d} + \varepsilon_d,
\qquad
\varepsilon_d\sim\mathcal N(0,R_{d}),
\qquad d=1,\dots,n_{\mathrm{exp}},
\]
with positive definite \(R_{d}\).

Define
\[
m_k(y):=\int \ell(y\mid\theta')\,q^{(k)}(d\theta'),
\qquad
m(y):=\int \ell(y\mid\theta')\,q^{(0)}(d\theta'),
\]
and
\[
g_k(\theta):=
-\mathbb{E}_{Y\sim \ell(\cdot\mid\theta)}[\log m_k(Y)],
\qquad
g(\theta):=
-\mathbb{E}_{Y\sim \ell(\cdot\mid\theta)}[\log m(Y)].
\]

Then \(g_k\to g\) locally uniformly on \(\mathbb R^{n_\theta}\), and there exists \(C>0\) such that
\[
|g_k(\theta)|\le C(1+\|\theta\|^2)
\qquad\text{for all }k,\theta.
\]
\end{lemma}

\begin{proposition}[Consistency in the Linear--Gaussian setting]
\label{prop:consistency_LG}
Consider a fixed design $(w, u)$ providing, for some observation matrix $H_{i, d}$ and  affine offset $b_{i,d}$, the linear--Gaussian observation models
\[
y_{i, d} = H_{i, d} \theta + b_{i, d} + \varepsilon_{i, d},
\qquad \varepsilon_{i, d} \sim \mathcal N(0,R_{i, d}),
\quad i=1,\dots,M, ~~\text{and}~~ d=1,\dots,n_{exp}
\]
with positive definite $R_{i, d}$. Let $J_{\mathrm{EIG}}(p_0)$ denote the exact expected information gain for this model under the parameter prior $p_0$, and let $J_{\mathrm{tilt}}(p_0)$ denote the Gaussian tilting objective constructed
from $p_0$.
\begin{enumerate}
\item
If the prior is Gaussian, $p_0=\mathcal N(m_0,\Sigma_0)$, then the Gaussian tilting
surrogate with $\theta_\mathrm{ref}=m_0$ is exact, and
\[
J_{\mathrm{tilt}}(p_0) = J_{\mathrm{EIG}}(p_0) = \tfrac{1}{2}\log\det(I + \Sigma_0 F_M),
\]
where \(F_M\) is the final accumulated Fisher information matrix defined
by~\eqref{eq:accumulator_definition}.
\item
Let $(p_k)_{k\ge1}$ be a sequence of priors with finite second moment such that
$p_k \to p_0:=\mathcal N(m_0,\Sigma_0)$ in $W_2$. Then, with \(\theta_{\mathrm{ref}}^{(k)}\) set to the mean of \(p_k\), we have,
for each fixed design, the pointwise convergence
\[
J_{\mathrm{tilt}}(p_k)
\xrightarrow[k\to\infty]{}
J_{\mathrm{EIG}}(p_0).
\]
\end{enumerate}
\end{proposition} 

\begin{proof}
We prove the two assertions sequentially.

\noindent\textit{1. Exactness in the Gaussian case.}
Assume that the prior is Gaussian, $p_0=\mathcal N(m_0,\Sigma_0)$.
In the linear--Gaussian model, the posterior distribution after $i$
observations remains Gaussian, with precision matrix
\[
\Sigma_i^{-1}
=
\Sigma_0^{-1}
+
\sum_{j=1}^{i}
\sum_{d=1}^{n_{exp}}
w_{j,d}\, H_{j,d}^\top R_{j,d}^{-1} H_{j,d}
=
\Sigma_0^{-1}+F_i,
\]
where $F_i$ is the accumulated Fisher information defined in \eqref{eq:accumulator_definition}.
In particular, the posterior covariance $\Sigma_i$ is deterministic and
independent of the realized data.

The exact expected information gain admits the closed--form expression
\begin{equation}\label{eq:eig_gauss}
J_{\mathrm{EIG}}(p_0)
=
\sum_{i=1}^M \frac12 \log \det\!\left(I + \Sigma_{i-1} F_i^\Delta \right),
\end{equation}
where $\Sigma_{i-1}$ denotes the posterior covariance after $i-1$ steps.

Consider now the Gaussian tilting surrogate. Since $p_0$ is centered at $m_0$ and
$\theta_{\mathrm{ref}}=m_0$, tilting by
\(
\phi_i(\theta)
=
\exp\!\left(-\tfrac12(\theta-m_0)^\top F_i(\theta-m_0)\right)
\)
yields a Gaussian distribution $q_i$ with precision matrix $\Sigma_0^{-1}+F_i$ and hence
covariance $\Sigma_i$. Therefore, although the deterministic surrogate $q_i$ does not track the
data-dependent mean of the true posterior, it reproduces its exact covariance.
Since in the linear--Gaussian setting the incremental information gain depends
only on the prior covariance, the contribution computed from $q_{i-1}$ coincides
with the exact contribution given by \eqref{eq:eig_gauss}. Summing over
$i=1,\dots,M$ yields
\[
J_{\mathrm{tilt}}(p_0)=J_{\mathrm{EIG}}(p_0).
\]

\medskip

\noindent\textit{2. Convergence under $W_2$.}
Let $(p_k)_{k\ge1}$ be a sequence of probability measures with finite second moments such that
\(
p_k \to p_0 := \mathcal N(m_0,\Sigma_0)
\) in the $2$--Wasserstein distance.
Denoting by $m_k=\mathbb E_{p_k}[\theta]$ the mean of $p_k$, we define the tilting factors
\[
\phi_i^{(k)}(\theta)
=
\exp\!\left(-\tfrac12(\theta-m_k)^\top F_i (\theta-m_k)\right) \quad \text{and} \quad \phi_i(\theta)
=
\exp\!\left(-\tfrac12(\theta-m_0)^\top F_i (\theta-m_0)\right),
\]
as well as the probability measures
\[
q_i^{(k)}(d\theta)
\propto
\phi_i^{(k)}(\theta)p_k(d\theta)\quad \text{and} \quad 
q_i^{(0)}(d\theta)
\propto
\phi_i(\theta)p_0(d\theta).
\]
From Lemma~\ref{lmm:weak_conv_qk}, we obtain that for any sequence of continuous functions $(f_k)$ converging locally uniformly to $f$ and satisfying a uniform quadratic bound, we have
\begin{equation}
\mathbb{E}_{q_i^{(k)}}[f_k(\theta)]
\;\longrightarrow\;
\mathbb{E}_{q_i^{(0)}}[f(\theta)]. \label{eq:convergence_qk}
\end{equation}
For each stage \(i\), let \(Y_i\) denote the observation block at step \(i\),
with conditional density
\[
\ell_i(y\mid\theta):=p(Y_i=y\mid\theta,w,u).
\]
To conclude, it remains to verify that the functions
\[
f_i(\theta)
:=
-\mathbb{E}_{Y_i\sim \ell_i(\cdot\mid\theta)}
[\log \ell_i(Y_i\mid\theta)]
\]
and
\[
g_{i,k}(\theta)
:=
-\mathbb{E}_{Y_i\sim \ell_i(\cdot\mid\theta)}
[\log m_{i,k}(Y_i)] \quad \text{where} \quad m_{i,k}(y)
=
\int \ell_i(y\mid\theta')\,q_{i-1}^{(k)}(d\theta')
\]


satisfy a uniform quadratic bound in \(\theta\), and that \(g_{i,k}\) converges locally uniformly to 
\[
g_i(\theta)
:=
-\mathbb{E}_{Y_i\sim \ell_i(\cdot\mid\theta)}
[\log m_i(Y_i)]
\qquad \text{where} \quad
m_i(y)
=
\int \ell_i(y\mid\theta')\,q_{i-1}^{(0)}(d\theta').
\]
In the linear--Gaussian observation model the log-likelihood
$\log \ell_i(y\mid\theta)$ grows at most quadratically in $(\theta,y)$, which directly yields the uniform quadratic bound on $f_i$. The desired properties on \(g_{i,k}\) follow from Lemma~\ref{lmm:behavior_denom_eig_gauss}.

We can thus apply \eqref{eq:convergence_qk} to the first term with the fixed function \(f_i\), and to the second term with \(g_{i,k}\to g_i\), to conclude that the \(i\)-th stage contribution converges. Summing over \(i=1,\dots,M\), we obtain
\[
\lim_{k\to\infty} J_{\mathrm{tilt}}(p_k)
=
J_{\mathrm{tilt}}(p_0).
\]
Combining this with Part~1 yields
\[
\lim_{k\to\infty} J_{\mathrm{tilt}}(p_k)
=
J_{\mathrm{EIG}}(p_0).
\]
\end{proof}

\end{document}
